\section{Lean formalization architecture}\label{sec:lean-plan}

\subsection{Integration target}

The proposed Lean namespace is
\path{LeanProofs.IPRA.NativeClauseFramework}.  It depends on two checked
libraries in the root workspace:

\begin{itemize}
\item \path{ConciseIPRA.NativeClausalDiagonal} and
\path{ConciseIPRA.NativeBiInterpretation};
\item \path{IPRADerivations.Core.Tactic} and later domain modules.
\end{itemize}

The implementation should be added to the existing
\path{Logic/PrimitiveRecursiveArithmetic/Derivations/Lean} Lake library rather
than maintained as a disconnected packet.  This allows every new application
to end in the same embedded \(\Prf\;\PRAAx\) calculus and reuse the checked
arithmetic stack.

\subsection{Module graph}

\begin{center}
\begin{tikzpicture}[node distance=9mm and 18mm, every node/.style={draw,rounded corners,align=center,inner sep=4pt,font=\small}, >=Latex]
\node (up) {upstream\\KPrf + bridges};
\node (bridge) [below=of up] {Bridge};
\node (program) [below left=of bridge] {Program};
\node (bool) [below right=of bridge] {Boolean};
\node (trace) [below=of bool] {Trace/Search};
\node (demand) [below=of program] {Demand};
\node (rel) [below=of trace] {Relational};
\node (domain) [below left=of rel] {Domains};
\node (apps) [below right=of rel] {Applications};
\draw[->] (up) -- (bridge);
\draw[->] (bridge) -- (program);
\draw[->] (bridge) -- (bool);
\draw[->] (bool) -- (trace);
\draw[->] (program) -- (demand);
\draw[->] (program) -- (rel);
\draw[->] (trace) -- (rel);
\draw[->] (demand) -- (domain);
\draw[->] (rel) -- (domain);
\draw[->] (domain) -- (apps);
\draw[->] (rel) -- (apps);
\end{tikzpicture}
\end{center}

The archive includes files with these names.  \path{Bridge.lean} is written as
candidate direct code using current upstream declarations.  The other files are
blueprints with explicit theorem signatures and \path{sorry}; they are not
presented as compiled modules.

\subsection{Bridge module}

The first module should prove only four small theorems:

\begin{enumerate}
\item \path{cutAsmList}: discharge a list of ordinary assumptions using
\path{KPrf.cut};
\item \path{cutOne}: one-output specialization of list cut;
\item \path{exportPRA}: compose \path{KPrf.toNPrf} with
\path{native_iff_pra};
\item \path{importPRA}: compose \path{pra_iff_native} with
\path{kprf_iff_nprf}.
\end{enumerate}

The intended declarations are:

\begin{lstlisting}[language={},caption={Core proof-exchange API}]
theorem cutAsmList
    (h : KPrf (Delta.map Cl.asm ++ Gamma) A)
    (hs : forall X in Delta, KPrf Gamma X) : KPrf Gamma A

theorem exportPRA (h : KPrf Gamma A) :
    Gamma.map (fun c => c.toN.toPRA) |- [PRAAx] A.toPRA

theorem importPRA (h : Gamma |- [PRAAx] phi) :
    KPrf ((Gamma.map NFm.ofPRA).map
      (fun X => Cl.asm (NFm.unQ X)))
      (NFm.unQ (NFm.ofPRA phi))
\end{lstlisting}

The exact notation in the artifact follows upstream syntax.  Because no Lean
binary was present, even these short declarations are labeled ``candidate''
until checked in the repository workspace.

\subsection{Dependent program datatype}

A program should record the earlier outputs available to each node.
One possible datatype is:

\begin{lstlisting}[language={},caption={Suggested dependent schedule}]
structure Node (n : Nat) (Gamma : List (Cl n))
    (available : List (NFm n)) where
  out : NFm n
  local : KPrf (available.map Cl.asm ++ Gamma) out

inductive Program (n : Nat) (Gamma : List (Cl n)) :
    List (NFm n) -> Type
  | nil : Program n Gamma []
  | snoc {outs} : Program n Gamma outs ->
      Node n Gamma outs -> Program n Gamma (node.out :: outs)
\end{lstlisting}

The order of \path{outs} can be newest-first to make extension constant-time.
The compiler returns a proof environment:

\begin{lstlisting}[language={}]
def Compiled (Gamma : List (Cl n)) (outs : List (NFm n)) : Prop :=
  forall A, A in outs -> KPrf Gamma A

theorem Program.compile (p : Program n Gamma outs) :
  Compiled Gamma outs
\end{lstlisting}

The proof is induction on \path{p}; the step is \path{cutAsmList}.  A terminal
method extracts the proof of the newest output.  This theorem should be the
first new kernel-checked milestone, because it validates the architecture with
minimal domain machinery.

\subsection{Sparse dependencies}

For large graphs, a node should not assume every previous output.  Store a list
of indices \path{deps : List (Fin available.length)} and define the dependency
formula list by lookup.  The compiler maps each index to the corresponding
compiled proof.  Correctness is still list cut.

An alternative is a named finite map.  Indices are preferable in the kernel
layer because membership and lookup are structurally simple; human-readable
names belong in generated metadata.

\subsection{Formula builders}

Application code should not manipulate raw de Bruijn indices everywhere.
Create typed builders for common forms:

\begin{itemize}
\item equations and Boolean equations;
\item bounded predicates compiled to characteristic terms;
\item tuple/list access equations;
\item guarded clauses;
\item trace-state invariants;
\item rational error inequalities;
\item relation and observer formulas.
\end{itemize}

Each builder should come with substitution lemmas and, where possible,
\path{simp} theorems.  The builders elaborate to ordinary \(\NFm\) or
\(\mathrm{Fm}\;\PTm\), so they do not enlarge the trusted kernel.

\subsection{Booleanization module}

There are two implementation strategies.

\paragraph{Reuse existing reflection.}
Audit \path{TeqReflect.lean} and related files for characteristic terms and
quantifier-free induction.  Add adapters to the \path{IPRADerivations} term
symbols and native formulas.

\paragraph{A small dedicated AST.}
Define a restricted Boolean formula syntax with equality, not, and, or, and
bounded folds.  Compile it to a PR characteristic term and to an ordinary
formula, then prove equivalence by structural induction.  This is easier to use
for generated verifiers and avoids proving reflection for arbitrary syntax.

The second route is recommended for the first milestone.  General reflection
can be integrated later.

\subsection{Trace module}

A direct ordinary-iPRA API can build on \path{IPRADerivations.ind}:

\begin{lstlisting}[language={}]
structure TraceSpec (k : Nat) where
  state : PTm (k + 1)
  inv : Fm PTm (k + 1)
  qf : inv.QF
  base : [] |- [PRAAx] inv.subst (inst zero)
  step : [inv] |- [PRAAx] inv.subst (repl (succ (var 0)))

theorem TraceSpec.closed (S : TraceSpec k) :
  [] |- [PRAAx] S.inv := ind S.inv S.qf S.base S.step
\end{lstlisting}

A higher-level builder constructs the state term from an initializer and step
PR symbol and constructs \path{base} and \path{step} from local equations.
Search, bounded folds, table construction, and demand passes are specializations.

\subsection{Demand module}

The pure metatheoretic layer can first be implemented over
\path{Vector Nat k} or \path{Fin k -> Nat} to prove the finite lattice theorems.
The object layer then chooses a concrete tuple code and PR functions.  Required
theorems are:

\begin{enumerate}
\item componentwise order and join laws;
\item rank strictness for a strict componentwise increase;
\item bounded closure termination;
\item fixed-point and leastness;
\item reverse DAG propagation correctness and minimality;
\item schedule lookup and replay correctness.
\end{enumerate}

Do not stop at the metatheoretic theorem.  The endpoint is a
\(\Prf\;\PRAAx\) trace theorem for the encoded closure algorithm.  The
metatheory guides the proof and tests the definitions.

\subsection{Relational module}

Define bounded clouds by a bound and Boolean predicate rather than by Mathlib
finite sets.  This representation maps directly to PR terms.  The module should
provide:

\begin{itemize}
\item reverse viability table;
\item least viable witness and least successor;
\item path coding and edge correctness;
\item exact observer invariance;
\item rational diameter check;
\item motion-tail Cauchy estimate;
\item gap-based exact recovery.
\end{itemize}

The general theorems should be proved once in ordinary iPRA and imported into
the native graph as needed.

\subsection{Algebra-node import}

A small adapter should turn an \path{IPRADerivations} theorem into a native node.
For a direct theorem \path{h}, define the node output to be
\path{NFm.unQ (NFm.ofPRA phi)} and the local proof to be
\path{importPRA h}.  The node's predecessor formulas are the translations of
\path{h}'s assumptions.

The reverse adapter exports a terminal node with \path{exportPRA}.  Round-trip
theorems can be used to rewrite the result to the human-facing formula rather
than leaving a translated formula in the public API.

\subsection{Suggested completion sequence}

\begin{longtable}{L{0.09\textwidth}L{0.42\textwidth}L{0.39\textwidth}}
\toprule
Step & Deliverable & Acceptance criterion \\
\midrule
1 & \path{Bridge.lean} & Compiles; no \path{sorry}; axiom audit no stronger than upstream metatheory.\\
2 & \path{Program.lean} & A two-node and a five-node graph compile to closed \(\KPrf\) proofs.\\
3 & Algebra exchange demo & Import one checked Heron-related identity, compose a native equality node, export to the same ordinary formula.\\
4 & Boolean formulas & Characteristic-term compiler for the restricted verifier language, with object-level equivalence.\\
5 & Trace and search & Closed derivations for a counter trace, least bounded search, and complete enumeration.\\
6 & Demand & Object-level finite monotone closure and a small cyclic example.\\
7 & Relational & Object-level three-layer viability and stable observer example.\\
8 & Domain library & Rational arithmetic, finite lists/tables, and dyadic errors.\\
9 & Packet migration & Convert one existing proof packet end to end, keeping a claim ledger.\\
10 & Automation & Generate node syntax and proof schedules from a declarative JSON or Lean DSL.\\
\bottomrule
\end{longtable}

\subsection{Testing and axiom audits}

Every checked theorem module should have an audit file using \path{#print axioms}
or the repository's established mechanism.  Every generated graph should have:

\begin{itemize}
\item a small executable reference example;
\item a test that deliberately corrupts one local certificate and is rejected;
\item a theorem statement snapshot to detect accidental weakening;
\item a count of \path{sorry}, expected to be zero in checked modules;
\item a comparison between the public ordinary formula and the exported
translation.
\end{itemize}

The proof kernel, not the tests, remains the authority.  Tests protect the code
generator and theorem statement.

\subsection{Avoiding common Lean traps}

\begin{enumerate}
\item Do not define the main theorem as a structure field and then count the
structure constructor as a proof.
\item Do not use a host theorem \path{eval spec = true} as the final endpoint
unless a proved reflection theorem constructs the object derivation.
\item Keep semantic interpretations and syntactic derivations in different
namespaces.
\item Normalize contexts before invoking bridge theorems; most failures are
list-map or \path{unQ} mismatches.
\item Prefer checked structural lemmas to rewriting under binders by hand.
\item Expose all conditional producer assumptions in the public theorem type.
\item Pin the upstream commit because the native calculus is under active
research and names may change.
\end{enumerate}
